% Options for packages loaded elsewhere
% Options for packages loaded elsewhere
\PassOptionsToPackage{unicode}{hyperref}
\PassOptionsToPackage{hyphens}{url}
\PassOptionsToPackage{dvipsnames,svgnames,x11names}{xcolor}
%
\documentclass[
  british,
]{article}
\usepackage{xcolor}
\usepackage{amsmath,amssymb}
\setcounter{secnumdepth}{5}
\usepackage{iftex}
\ifPDFTeX
  \usepackage[T1]{fontenc}
  \usepackage[utf8]{inputenc}
  \usepackage{textcomp} % provide euro and other symbols
\else % if luatex or xetex
  \usepackage{unicode-math} % this also loads fontspec
  \defaultfontfeatures{Scale=MatchLowercase}
  \defaultfontfeatures[\rmfamily]{Ligatures=TeX,Scale=1}
\fi
\usepackage{lmodern}
\ifPDFTeX\else
  % xetex/luatex font selection
\fi
% Use upquote if available, for straight quotes in verbatim environments
\IfFileExists{upquote.sty}{\usepackage{upquote}}{}
\IfFileExists{microtype.sty}{% use microtype if available
  \usepackage[]{microtype}
  \UseMicrotypeSet[protrusion]{basicmath} % disable protrusion for tt fonts
}{}
\makeatletter
\@ifundefined{KOMAClassName}{% if non-KOMA class
  \IfFileExists{parskip.sty}{%
    \usepackage{parskip}
  }{% else
    \setlength{\parindent}{0pt}
    \setlength{\parskip}{6pt plus 2pt minus 1pt}}
}{% if KOMA class
  \KOMAoptions{parskip=half}}
\makeatother
% Make \paragraph and \subparagraph free-standing
\makeatletter
\ifx\paragraph\undefined\else
  \let\oldparagraph\paragraph
  \renewcommand{\paragraph}{
    \@ifstar
      \xxxParagraphStar
      \xxxParagraphNoStar
  }
  \newcommand{\xxxParagraphStar}[1]{\oldparagraph*{#1}\mbox{}}
  \newcommand{\xxxParagraphNoStar}[1]{\oldparagraph{#1}\mbox{}}
\fi
\ifx\subparagraph\undefined\else
  \let\oldsubparagraph\subparagraph
  \renewcommand{\subparagraph}{
    \@ifstar
      \xxxSubParagraphStar
      \xxxSubParagraphNoStar
  }
  \newcommand{\xxxSubParagraphStar}[1]{\oldsubparagraph*{#1}\mbox{}}
  \newcommand{\xxxSubParagraphNoStar}[1]{\oldsubparagraph{#1}\mbox{}}
\fi
\makeatother


\usepackage{longtable,booktabs,array}
\newcounter{none} % for unnumbered tables
\usepackage{calc} % for calculating minipage widths
% Correct order of tables after \paragraph or \subparagraph
\usepackage{etoolbox}
\makeatletter
\patchcmd\longtable{\par}{\if@noskipsec\mbox{}\fi\par}{}{}
\makeatother
% Allow footnotes in longtable head/foot
\IfFileExists{footnotehyper.sty}{\usepackage{footnotehyper}}{\usepackage{footnote}}
\makesavenoteenv{longtable}
\usepackage{graphicx}
\makeatletter
\newsavebox\pandoc@box
\newcommand*\pandocbounded[1]{% scales image to fit in text height/width
  \sbox\pandoc@box{#1}%
  \Gscale@div\@tempa{\textheight}{\dimexpr\ht\pandoc@box+\dp\pandoc@box\relax}%
  \Gscale@div\@tempb{\linewidth}{\wd\pandoc@box}%
  \ifdim\@tempb\p@<\@tempa\p@\let\@tempa\@tempb\fi% select the smaller of both
  \ifdim\@tempa\p@<\p@\scalebox{\@tempa}{\usebox\pandoc@box}%
  \else\usebox{\pandoc@box}%
  \fi%
}
% Set default figure placement to htbp
\def\fps@figure{htbp}
\makeatother



\ifLuaTeX
\usepackage[bidi=basic,shorthands=off]{babel}
\else
\usepackage[bidi=default,shorthands=off]{babel}
\fi
\ifLuaTeX
  \usepackage{selnolig} % disable illegal ligatures
\fi


\setlength{\emergencystretch}{3em} % prevent overfull lines

\providecommand{\tightlist}{%
  \setlength{\itemsep}{0pt}\setlength{\parskip}{0pt}}



 


% Get Math Fonts
\usepackage{fontspec}
\usepackage{unicode-math}
\setmathfont{Latin Modern Math}
\setmathfont[range={\mathscr,\mathbfscr,\mathbb}]{XITSMath-Regular}
[    Extension = .otf,
      BoldFont = XITSMath-Bold
]
\makeatletter
\@ifpackageloaded{caption}{}{\usepackage{caption}}
\AtBeginDocument{%
\ifdefined\contentsname
  \renewcommand*\contentsname{Table of contents}
\else
  \newcommand\contentsname{Table of contents}
\fi
\ifdefined\listfigurename
  \renewcommand*\listfigurename{List of Figures}
\else
  \newcommand\listfigurename{List of Figures}
\fi
\ifdefined\listtablename
  \renewcommand*\listtablename{List of Tables}
\else
  \newcommand\listtablename{List of Tables}
\fi
\ifdefined\figurename
  \renewcommand*\figurename{Figure}
\else
  \newcommand\figurename{Figure}
\fi
\ifdefined\tablename
  \renewcommand*\tablename{Table}
\else
  \newcommand\tablename{Table}
\fi
}
\@ifpackageloaded{float}{}{\usepackage{float}}
\floatstyle{ruled}
\@ifundefined{c@chapter}{\newfloat{codelisting}{h}{lop}}{\newfloat{codelisting}{h}{lop}[chapter]}
\floatname{codelisting}{Listing}
\newcommand*\listoflistings{\listof{codelisting}{List of Listings}}
\usepackage{amsthm}
\theoremstyle{definition}
\newtheorem{definition}{Definition}[section]
\theoremstyle{plain}
\newtheorem{proposition}{Proposition}[section]
\theoremstyle{remark}
\AtBeginDocument{\renewcommand*{\proofname}{Proof}}
\newtheorem*{remark}{Remark}
\newtheorem*{solution}{Solution}
\newtheorem{refremark}{Remark}[section]
\newtheorem{refsolution}{Solution}[section]
\makeatother
\makeatletter
\makeatother
\makeatletter
\@ifpackageloaded{caption}{}{\usepackage{caption}}
\@ifpackageloaded{subcaption}{}{\usepackage{subcaption}}
\makeatother

\newcommand{\Measure}{{\mu}}
\newcommand{\ProjectionMap}[1][\vphantom{}]{{\pi}_{#1}}
\newcommand{\SigmaAlgebra}[1]{{\mathscr{#1}}}
\newcommand{\Expectation}[1][\Measure]{{\mathbb{E}}_{#1}}
\NewDocumentCommand{\ProductMeasure}{O{\Measure} O{\nu} }{{#1}}

\usepackage{bookmark}
\IfFileExists{xurl.sty}{\usepackage{xurl}}{} % add URL line breaks if available
\urlstyle{same}
% fallback for those not using the hyperref driver hyperxmp:
\makeatletter
\@ifundefined{xmpquote}{\newcommand{\xmpquote}[1]{#1}}{}
\makeatother
\hypersetup{
  pdftitle={Joinings},
  pdflang={en-GB},
  colorlinks=true,
  linkcolor={blue},
  filecolor={Maroon},
  citecolor={Blue},
  urlcolor={Blue},
  pdfcreator={LaTeX via pandoc}}


\title{Joinings}
\author{}
\date{2026-09-26}
\begin{document}
\maketitle


We want to build our understanding of joinings as they help us create a
generalised definition of a measure on product spaces that are behave as
we would expect for a multi-dimensional measure. This is important when
we start looking at \(n\)-dimensional Erdős cubes. In building this
understanding, we used the product measure as an example and
investigated what property we are interested in that we'd like to see in
joinings.

\begin{definition}[]\protect\hypertarget{def-ProductMeasure}{}\label{def-ProductMeasure}

The \emph{product measure} on
\((X\times Y,\SigmaAlgebra{B}\otimes\SigmaAlgebra{C})\) is
\(\ProductMeasure\) where
\[(\ProductMeasure)(B\times C)=\Measure(B)\nu(C)\] for
\(B\in\SigmaAlgebra{B},C\in\SigmaAlgebra{C}\).

The product measure is \(T\times S\)-invariant, i.e., for all
\(D\in\SigmaAlgebra{B}\otimes\SigmaAlgebra{C}\),
\[(\ProductMeasure)(D)=(\ProductMeasure)((T\times S)^{-1}D).\]

Note, if \(D=B\times C\), then
\[(T\times S)^{-1}(B\times C)=T^{-1}B\times S^{-1}C.\]

\end{definition}

The main property we are interested in relate to how we can relate the
measure of a product space to the measures of the component spaces. This
leads us to the definition of marginals:

\begin{definition}[]\protect\hypertarget{def-Marginals}{}\label{def-Marginals}

Let \(\ProjectionMap[X]:X\times Y\rightarrow X\) and
\(\ProjectionMap[Y]:X\times Y\rightarrow Y\) be coordinate projections
of a probability space
\((X\times Y,\SigmaAlgebra{B}\times\SigmaAlgebra{C},\eta)\). We call the
coordinate-projected measures \(\ProjectionMap[X](\eta)\) and
\(\ProjectionMap[Y](\eta)\) the \emph{marginals} of \(\eta\).

These have the property that
\[(\ProjectionMap[X](\eta))(B)=\eta(\ProjectionMap[X]^{-1}(B))\] and
\[(\ProjectionMap[Y](\eta))(C)=\eta(\ProjectionMap[Y]^{-1}(C)).\]

\end{definition}

When looking at the product measure \(\eta=\ProductMeasure\) as our
example, we get the marginal \(\Measure\) as \begin{align*}
    (\ProjectionMap[X](\ProductMeasure))(B)&=(\ProductMeasure)(\ProjectionMap[X]^{-1}(B))
    \\&=(\ProductMeasure)(B\times Y)
    \\&=\Measure(B)\nu(Y)
    \\&=\Measure(B),
\end{align*} so \(\ProjectionMap[X](\ProductMeasure)=\Measure\) and we
get \(\ProjectionMap[Y](\ProductMeasure)=\nu\) similarly. Hence, when
looking at measures on product spaces, we want to conserve this
property. Thus, we arrive at the definition of joinings:

\begin{definition}[]\protect\hypertarget{def-Joinings}{}\label{def-Joinings}

A \emph{joining} of \((X,\SigmaAlgebra{B},\Measure,T)\) and
\((Y,\SigmaAlgebra{C},\nu,S)\) is a measure \(\eta\) on
\((X\times Y,\SigmaAlgebra{B}\otimes\SigmaAlgebra{C})\) satisfying:

\begin{enumerate}
\def\labelenumi{\arabic{enumi}.}
\tightlist
\item
  \(\eta\) is \(T\times S\)-invariant,
\item
  and, \(\eta\) has marginals \(\Measure\) and \(\nu\).
\end{enumerate}

A \emph{self-joining} of \((X,\SigmaAlgebra{B},\Measure,T)\) is a
joining of \((X,\SigmaAlgebra{B},\Measure,T)\) and itself.

\end{definition}

The product measure is always a joining and it may sometimes be the only
joining for some product spaces. For our research, we're interested in
the product space of \((X,\SigmaAlgebra{B},\Measure)\) with itself, so
we will be focusing on self-joinings. We move on to this special case of
self-joinings:

\begin{proposition}[]\protect\hypertarget{prp-DiagonalJoining}{}\label{prp-DiagonalJoining}

Let \[\triangle=\{(x,x):x\in X\}\subseteq X\times X.\] Then there is a
self-joining \(\eta\) of \((X,\SigmaAlgebra{B},\Measure,T)\) with
\(\eta(\triangle)=1\).

\end{proposition}

For the product measure, we only get \((\ProductMeasure)(\triangle)>0\)
if \(\Measure(\{a\})>0\), for atoms \(a\in X\), and we get
\((\ProductMeasure)(\triangle)=1\) if and only if \(X=\{*\}\), i.e.,
\(X\) only has a single point.

To show that we will there exists a self-joining of
\((X,\SigmaAlgebra{B},\Measure,T)\), \(\eta\), such that
\(\eta(\triangle)=1\), we define \(\phi:X\rightarrow X\times X\) by
\(\phi(x)=(x,x)\) and define the push-forward measure of \(\Measure\)
under \(\phi\) as \begin{align*}
    \phi\Measure(B\times C) &=\Measure(\phi^{-1}(B\times C))
    \\&=\Measure(\{x\in X: x\in B,x\in C \})
    \\&=\Measure(B\cap C).
\end{align*} As \(\phi^{-1}(\triangle)=X\) and \(\Measure(X)=1\), we
arrive at \(\phi\Measure(\triangle)=1\) and by setting
\(\eta=\phi\Measure\), we complete the proof.

\section{Relatively Independent
Joinings}\label{relatively-independent-joinings}

Before we can look at the measures required for Erd"\{o\}s cubes, we
also set out to understand relatively independent joinings.

\begin{definition}[]\protect\hypertarget{def-RelativelyIndependentJoining}{}\label{def-RelativelyIndependentJoining}

A \emph{relatively independent joining} of
\((X,\SigmaAlgebra{B},\Measure,T)\) is the self joining, \(\nu\), of
\((X,\SigmaAlgebra{C},\Measure,T)\) where \(\SigmaAlgebra{C}\) is the
\(\sigma\)-subalgebra of \(T\)-invariant subsets in
\(\SigmaAlgebra{B}\).

\end{definition}

We construct an example of a relatively independent joining.

Fix \((X,\SigmaAlgebra{B},\Measure,T)\) and a \(\sigma\)-subalgebra
\(\SigmaAlgebra{C}\subseteq\SigmaAlgebra{B}\) such that
\(\SigmaAlgebra{C}\) is \(T\)-invariant. This gives us the following
properties for all \(f\in\text{L}^2(X,\SigmaAlgebra{B},\Measure)\):

\begin{enumerate}
\def\labelenumi{\arabic{enumi}.}
\tightlist
\item
  \(\Expectation(f\mid\SigmaAlgebra{C})\) exists.
\item
  \(\Expectation(Tf\mid\SigmaAlgebra{C})\) exists as \(Tf=f\circ T\).
\item
  \(T\ \Expectation(f\mid\SigmaAlgebra{C})=\Expectation(Tf\mid\SigmaAlgebra{C})\).
\end{enumerate}

This allows us to use \(\SigmaAlgebra{C}\) to define a self-joining
\(\nu\) such that
\[\nu(A\times B)=\int\Expectation(\mathbb{1}_A\mid\SigmaAlgebra{C})\cdot\Expectation(\mathbb{1}_B\mid\SigmaAlgebra{C})\ \mathrm{d}\Measure.\]
We know this is a probability measure as
\(\nu(X\times X)=\int\mathbb{1}_X\cdot\mathbb{1}_X\ \mathrm{d}\Measure=\Measure(X)=1\).
Let \(\ProjectionMap[1],\ProjectionMap[2]:X\times X\rightarrow X\) be
our projections of the first and second coordinates, respectively.
Fixing \(B\in\SigmaAlgebra{B}\), \begin{align*}
    \nu(\ProjectionMap[1]^{-1}(B))=\nu(B\times X)
    &=\int\Expectation(\mathbb{1}_B\mid\SigmaAlgebra{C})\cdot\Expectation(\mathbb{1}_X\mid\SigmaAlgebra{C})\ \mathrm{d}\Measure
    \\&=\int\Expectation(\mathbb{1}_B\mid\SigmaAlgebra{C})\ \mathrm{d}\Measure
    \\&=\Measure(B),
\end{align*} as \(\Expectation(\mathbb{1}_X\mid\SigmaAlgebra{C})=1\).
Thus, we have shown that the marginal \(\ProjectionMap[1](\nu)\) is
\(\Measure\) and we also have \(\ProjectionMap[2](\nu)=\Measure\)
similarly. The proof that \(\nu\) is \(T\times T\)-invariant is also
trivial.

We also proved: If \((X,\SigmaAlgebra{B},\Measure,T)\) is ergodic, then
the relatively independent joining of
\((X,\SigmaAlgebra{B},\Measure,T)\) is the product measure
\(\Measure\times\Measure\), and that this may not be true for the
relatively independent joining of
\((X\times X,\SigmaAlgebra{B}\times\SigmaAlgebra{B},\Measure\times\Measure,T\times T)\)
as
\((X\times X,\SigmaAlgebra{B}\times\SigmaAlgebra{B},\Measure\times\Measure,T\times T)\)
may not be an ergodic system.




\end{document}
